%% showcase.tex -- what the grown margin does besides growing: it reads the text.
%% LuaLaTeX, two passes (the option aging reads the page count of the first):
%%   lualatex showcase.tex && lualatex showcase.tex
\documentclass[11pt,a4paper]{article}
\renewcommand\rmdefault{cmr}\renewcommand\sfdefault{cmss}\renewcommand\ttdefault{cmtt}
\usepackage[T1]{fontenc}
\usepackage[a4paper,left=68mm,right=22mm,top=32mm,bottom=30mm,headsep=8mm]{geometry}
\usepackage{fancyhdr}
\usepackage[grown,sections,aging]{florilegium}

\pagestyle{fancy}\fancyhf{}
\fancyhead[R]{\small\textsc{A year in the margin}}
\fancyfoot[R]{\small\thepage}
\renewcommand\headrulewidth{0pt}
\newcommand\code[1]{\texttt{\small #1}}

\begin{document}
\section{Spring}

\florilegiuminitial[5]{E}{very} page of this document grows its own garden, but the
gardens here also read the text they surround. This first page is spring: the plant is
young, its leaves are small and pale, and it is mostly in bud; the last page is autumn.
The document loads the package with three options, and the sections that follow show what
each of them does, along with a few things that need no option at all:
\begin{center}\code{\textbackslash usepackage[grown,sections,aging]\{florilegium\}}\end{center}

The letter that opens this paragraph is one of them. It was set with
\code{\textbackslash florilegiuminitial[5]\{E\}}, which reads the outline of the letter from an OpenType font (Latin
Modern Bold, the OpenType twin of Computer Modern, unless another is named), paints it
like a petal, and sets it in a tile five lines tall. A vine crosses the tile and passes
over and under the strokes of the letter in turn; where it passes over, a thin halo of
paper parts it from the letter. Leaves go where they fit and small flowers into the
largest gaps, which for a B or an O are its counters. The letter is also typeset in the
text, invisibly, so the paragraph still copies and searches as \emph{Every}.

\subsection{Flowers at the headings}

With the option \code{sections}, every numbered section and subsection gets a flower in the
margin, beside its heading: a large one for a section, a spray of small ones for a
subsection, painted in the copper of the accents, the way headings were once picked out in
red. The heading carries an invisible mark. When the page is shipped out, the
garden is drawn and the page box is searched for these marks, so the garden knows the
height of each heading on the same page in the same run. Where the margin is already
taken, by the large flower at the top of the page for example, the flower is made
smaller or left out.

\subsection{A flower where the text asks for one}

The same mark can be placed by hand. \code{\textbackslash florilegiumbloom} puts a flower
in the margin beside the line in which it stands, like this
one\florilegiumbloom, and \code{\textbackslash florilegiumbloom[lily]} asks for any of the
motifs the garden is made of: \code{lily}, \code{blossoms}, \code{cyme}, \code{umbel},
\code{bud}, or a \code{butterfly}. It costs nothing on pages that have no marks and does
not move the text.

\section{Summer}

By the middle of the document the garden is in full bloom: open flowers, star lilies,
forget-me-nots curling at their tips, umbels. The page count needed to know where the
middle is comes from the previous run, which is why the option \code{aging} asks for two;
on the first run every page is summer, and LaTeX says so with a request to rerun.

The seasons change what grows, not where. In spring sprays of closed buds take the place of
most open flowers, the two shorter stems end in buds instead of a lily and a flower, the
leaves are smaller and paler, and there are more tendrils. Summer is the garden as it
always was, which is also what a document without the option gets.

\subsection{Visitors}

Most pages have visitors: one or two butterflies and bees in the open spaces, a snail
crawling up one of the main stems, and on about one page in five a bird perched on a twig
of its own. They are drawn from their own random stream, so a page without them is
exactly the page it would have been, and the option \code{novisitors} leaves them all out.
In the growth animation the butterflies and bees fly in, the bird glides onto its twig
and the snail crawls up its stem until the last frame.

\subsection{Watching it grow}

Every garden can be watched as it grows. \code{make grow-7} records the growth of seed 7
once, while the page is laid out: every stem is also drawn as shorter prefixes of itself,
and every motif at a few smaller sizes with the same random draws. Each frame is then a
page that writes the pieces whose time has come, in their original painting order, so the
last frame is exactly the finished page. A short script turns the frames into a GIF and an
MP4; \code{EXPLAIN=1} adds a dashed circle for each free space the loop picks, just before
a branch grows to it. The stems rise first and the leaves open as their tips pass; then
the large holes fill, slowly, and the small ones last, quickly.

\subsection{One seed, one garden}

Nothing changes from one run to the next. A page's seed comes from its page number, or
from \code{\textbackslash def\textbackslash florilegiumseed\{7\}} for all of them, and
every decision is drawn from random streams derived from it: the layout from one, each
primitive from a fork of another, and the visitors and the season each from their own.
A page without visitors, or without an age, is therefore exactly the page it would have
been without those features. The same source always gives the same PDF.

\subsection{Colors}

The whole palette comes from one color, called \code{ink}. Defining it with
\code{\textbackslash definecolor\{ink\}\{HTML\}\{5B2A1F\}} paints the garden in sepia, \code{2F6B4F} in green, \code{8E1B3A} in crimson: the ramp from
the darkest stroke to the palest wash is computed from that one hue, so the light end of
it never turns gray. A second color, \code{accent}, is the copper of the hairlines, the
sparkles and the flowers' eyes, and in autumn of some of the leaves.

\section{Autumn}

On the last page the garden goes to seed. Fewer flowers open; the shorter stems end in
seed heads, dandelion clocks whose seeds drift away from them; loose petals have fallen
into the space below the flowers; and some of the leaves have turned copper.
\code{\textbackslash florilegiumage\{0.8\}} sets the age by hand instead, from 0, spring,
through 0.5, summer, to 1, autumn, for all the pages that follow.

The text ends here, early on its page, and \code{\textbackslash florilegiumtextend} tells
the garden so: it grows into the space below these lines as well.
\florilegiumtextend
\end{document}
